\begin{abstract}
UTXO 接收方若须等待公开执行后才能花费，每一跳支付都会增加确认等待。Next 让接收方在一轮认证后继续付款。核心思想是：随输出携带的证书已指明缺失来源的责任承担者，即其直接发行方。消费组织成员锁定输入，并为全部认证输出预留额度；接收方核验法定人数证书（TXCer）后，可以请求后继付款，后继仍须满足自身的准入与额度条件。若后继先执行，公开执行原子地消费认证输出，并登记有资金支持的直接发行方义务。来源执行则履行该义务；否则，入块且有效的到期后公开决定从发行方备付金中赔付缺口，迟到来源偿还备付，始终不执行的来源则由发行方承担损失。授权历史修复在原区块身份下将扣款引用写入消费输入，无需重新执行支付，经济闭合也无需等待历史修复。在每个四成员组织及四验证者委员会至多有一名静态拜占庭参与者的模型下，我们证明消费唯一性、每份证书未结责任的资金覆盖及 CAL 守恒；对于有限、无环、来源闭合、无冲突且每笔均成功执行一次的支付集合，终态 CAL 持有量与执行次序无关。在单机、成员同步提交、钱包 NoSync 的设置下，100 跳链到达最后一笔认证到账的中位时间为 4.707\,s，逐跳等待模式为 64.479\,s。
\end{abstract}

\begin{IEEEkeywords}
区块链，认证支付，抵押，拜占庭容错，
变色龙哈希。
\end{IEEEkeywords}
